\clearpage
\section{Accuracy \& Loss Plots}
\label{app:loss_acc_curve}

\subsection{Accuracy Plots}
\begin{figure*}[h]
  \centering
  \begin{minipage}{.48\textwidth}
    \centering
    \includegraphics[width=\linewidth]{Pics and figs/lstm bin acc.png}
    \caption*{Binary Classification }
    \label{fig:binary_accuracy}
  \end{minipage}\hfill \hfill % Add horizontal space
  \begin{minipage}{.48\textwidth}
    \centering
    \includegraphics[width=\linewidth]{Pics and figs/lstm multi acc.png}
    \caption*{Multiclass Classification }
    \label{fig:multiclass_accuracy}
  \end{minipage}
\caption{Accuracy Plots of Deep Learning Models}  
\end{figure*}


In binary classification LSTM models exhibit consistent accuracy gains of approximately 1.5\% to 2.5\%, while multiclass LSTM models show varying improvements from 3\% to 6\%. Models with attention mechanisms showcase a 1\% to 2.5\% increase in accuracy. CNN with LSTM models, in both binary and multiclass scenarios, demonstrate accuracy improvements ranging from 6\% to 10\% over the 10 epochs, highlighting their ability to capture nuanced patterns and adapt effectively.
\vspace{0.5cm}

\begin{figure*}[h]
  \centering
  \begin{minipage}{.48\textwidth}
    \centering
    \includegraphics[width=\linewidth]{Pics and figs/trans bin acc.png}
    \caption*{Binary Classification }
    \label{fig:binary_accuracy}
  \end{minipage}\hfill \hfill % Add horizontal space
  \begin{minipage}{.48\textwidth}
    \centering
    \includegraphics[width=\linewidth]{Pics and figs/trans multi acc.png}
    \caption*{Multiclass Classification }
    \label{fig:multiclass_accuracy}
  \end{minipage}
\caption{Accuracy Plots of Transformer Models}  
\end{figure*}


The XLM-RoBERTa model for binary classification consistently performs, with accuracy gains of approximately 1.5\% to 2\% over the 5 epochs. BanglaBERT, in binary classification, demonstrates accuracy improvements ranging from 6\% to 8\%, showcasing its robust performance. In multiclass scenarios, XLM-RoBERTa achieves accuracy gains of around 1\% to 2\%, while BanglaBERT shows more substantial improvements, ranging from 9\% to 11\%, emphasizing its effectiveness in capturing intricate patterns across various classes.

\vspace{0.5cm}
\clearpage
\begin{figure*}[h]
  \centering
  \begin{minipage}{.48\textwidth}
    \centering
    \includegraphics[width=\linewidth]{Pics and figs/gcn bin acc.png}
    \caption*{Binary Classification }
    \label{fig:binary_accuracy}
  \end{minipage}\hfill \hfill % Add horizontal space
  \begin{minipage}{.48\textwidth}
    \centering
    \includegraphics[width=\linewidth]{Pics and figs/gcn multi accuracy.png}
    \caption*{Multiclass Classification }
    \label{fig:multiclass_accuracy}
  \end{minipage}
\caption{Accuracy Plots of BanglaBERTGCN}  
\end{figure*}


In binary classification, the BanglaBERTGCN model exhibits consistent improvements in accuracy, showing significant gains from the first to the second epoch and maintaining a strong upward trend into the third epoch. The multiclass classification results follow a similar pattern, with noticeable percentage improvements from the first to the second epoch and continued advancement into the third epoch, highlighting the model's effectiveness and robust learning capabilities.

\subsection{Loss Plots}
\begin{figure*}[h]
  \centering
  \begin{minipage}{.48\textwidth}
    \centering
    \includegraphics[width=\linewidth]{Pics and figs/lstm bin loss.png}
    \caption*{Binary Classification }
    \label{fig:binary_accuracy}
  \end{minipage}\hfill \hfill % Add horizontal space
  \begin{minipage}{.48\textwidth}
    \centering
    \includegraphics[width=\linewidth]{Pics and figs/lstm multi loss.png}
    \caption*{Multiclass Classification }
    \label{fig:multiclass_accuracy}
  \end{minipage}
\caption{Loss Plots of Deep Learning Models}  
\end{figure*}

In binary classification, the LSTM models display steady reductions in loss, with percentage improvements ranging from approximately 4\% to 9\% across the 10 epochs. Models with attention mechanisms demonstrate more substantial decreases, showing percentage improvements from around 32\% to 55\%. The CNN with LSTM models, in both binary and multiclass settings, exhibit consistent declines in loss, with improvements ranging from approximately 14\% to 27\%.

\clearpage
\begin{figure*}[h]
  \centering
  \begin{minipage}{.48\textwidth}
    \centering
    \includegraphics[width=\linewidth]{Pics and figs/trans bin loss.png}
    \caption*{Binary Classification }
    \label{fig:binary_accuracy}
  \end{minipage}\hfill \hfill % Add horizontal space
  \begin{minipage}{.48\textwidth}
    \centering
    \includegraphics[width=\linewidth]{Pics and figs/trans multi loss.png}
    \caption*{Multiclass Classification }
    \label{fig:multiclass_accuracy}
  \end{minipage}
\caption{Loss Plots of Transformer Models}  
\end{figure*}
In binary classification, XLM-RoBERTa and BANG-BERT models consistently reduce losses, with improvements of approximately 8\% to 45\% and 16\% to 33\%, respectively. For multiclass classification, both models demonstrate reductions in loss, with XLM-RoBERTa showing improvements from around 38\% to 66\%, and BANG-BERT ranging from approximately 56\% to 69\%. This indicates effective loss minimization across epochs for both transformer models in binary and multiclass scenarios.

\vspace{0.5cm}
\begin{figure*}[h]
  \centering
  \begin{minipage}{.48\textwidth}
    \centering
    \includegraphics[width=\linewidth]{Pics and figs/gcn bin loss.png}
    \caption*{Binary Classification }
    \label{fig:binary_accuracy}
  \end{minipage}\hfill \hfill % Add horizontal space
  \begin{minipage}{.48\textwidth}
    \centering
    \includegraphics[width=\linewidth]{Pics and figs/gcn multi loss.png}
    \caption*{Multiclass Classification }
    \label{fig:multiclass_accuracy}
  \end{minipage}
\caption{Loss Plots of BanglaBERTGCN}  
\end{figure*}

In binary classification, the loss exhibited an increase of approximately 24.84\% over the 10 epochs. Similarly, in multiclass classification, the loss grew by about 15.97\%. These trends highlight potential challenges in convergence and may prompt further investigation into optimizing the model's training dynamics.